%! Choosing a Sample — Unit 1, Lesson 1.3 — Exit Ticket
\documentclass[10pt]{article}
\usepackage{saar-article}
\usepackage{saar-boxes}
\newcommand{\TallMath}[1]{$\displaystyle #1\rule[-1.4em]{0pt}{3.2em}$}

\begin{document}

\pageheader{Unit 1, Lesson 1.3}{Exit Ticket}
% NAMESTRIP: no \namedateperiod here. The student writes their name once, on the
% cover this component is stapled behind. See references/conventions.md.

\vspace{0.15in}

\begin{notesbox}{Exit Ticket --- No Notes}
\small
Bayside Middle School has $750$ students in three grades. The principal wants to
survey $30$ of them about the lunch schedule.

\begin{enumerate}[leftmargin=1.6em, itemsep=9pt, topsep=3pt]

  \item Name the sampling technique in each plan. Choose from \emph{simple
        random}, \emph{stratified}, \emph{systematic}, \emph{cluster}, and
        \emph{convenience}.

        \begin{center}
        \renewcommand{\arraystretch}{1.45}
        \begin{tabularx}{0.96\linewidth}{@{} X p{2.9cm} @{}}
        \toprule
        \textbf{The plan} & \textbf{Technique} \\
        \midrule
        Draw $2$ of the school's $30$ homeroom numbers at random and survey every
        student in those two rooms. & \blank{2.5cm} \\
        Take a random sample of $10$ sixth graders, $10$ seventh graders, and
        $10$ eighth graders. & \blank{2.5cm} \\
        Ask the first $30$ students who come through the cafeteria door.
          & \blank{2.5cm} \\
        \bottomrule
        \end{tabularx}
        \end{center}

  \item The principal instead numbers all $750$ students and takes a systematic
        sample of $30$.

        \textbf{(a)} Find the interval $k$.

        \begin{work}
          k &= 750 \div 30 \\
            &= 25
        \end{work}

        \textbf{(b)} The random start lands on student $6$. Write the next two
        students chosen. \blank{2.2cm} \qquad \blank{2.2cm}

  \item Sixth graders and eighth graders have very different lunch periods, so
        they are likely to answer differently.

        \textbf{(a)} Which technique \emph{guarantees} all three grades are
        represented in proportion? \blank{4cm}

        \textbf{(b)} A simple random sample of $30$ is completely fair. Explain
        why it still does not \emph{guarantee} that result.
        \writelines{2}

\end{enumerate}
\end{notesbox}

\end{document}
